\section{Conclusion}\label{sec_conclusion}

We introduced in this paper new time-scale-frequency transforms that can adapt themselves to the analysed signal, through the frequency domain and time domain focus functions $\sigma^\nu$ and $\sigma^\tau$. These adapt dynamically the scale/bandwidth of analysis windows as a function of frequency or time, leading to non-linear transforms. Under suitable assumptions on focus functions, we could prove first important results on the transforms such as the well-definedness on $L^2$, and norm controls similar to the one obtained in the linear case.

We stress that due to the non-linearity of the transform, those controls cannot be understood in terms of operator norms. Exploring the ways to define norms on general map on Banach spaces (in our case $L^2$) would be an interesting perspective. A first way could be to follow~\cite{wei2020development}, where generalised operator norms are introduced, which cover the non-linear case. For $F:L^2(\R)\rightarrow L^2(\R^2)$, define
\begin{equation}
    \rho(F) := \max\left\{   \sup_{f\not=0}\frac{\|F(f)\|_{L^2}}{\|f\|_{L^2}}, \|F(0)\|_{L^2}  \right\} \in [0,+\infty]\ .
\end{equation}
%Such way to define a norm insist on the bounded aspect of the application. 
It is quite easy to see from theorem~\ref{theo_bounding_time_marseillan_transform} that the time-focus transform $M^\tau$ satisfies
\begin{equation}
    \rho(M^\tau) \leqslant \|h\|_{L^\infty}|\supp(h)|^{1/2}\ ,
\end{equation}
and has therefore a finite norm in the above sense. However at this point the control from Theorem~\ref{theorem_norm_control_frequency_focus} does not allow us to conclude a similar control for $\rho(M^\nu)$.

An alternative is to define the norm of $F$ as its Lipschitz constant (since continuity and boundedness coincide in the case of linear operators), \ie\, the existence of a constant $C(F)$ such that
\begin{equation}\label{eq_lipchitz_constant}
\forall f_1,f_2\in L^2\ ,\quad \|F(f_1) - F(f_2) \|_{L^2(\R^2)} \leqslant C(F)\|f_1-f_2\|_{L^2(\R)}\ .
\end{equation}
Proving the existence of such Lipschitz constants for $M^\nu$ and  $M^\tau$ would ensure their uniform continuity. We saw in Lemma~\ref{lemma_density_result} that the time-focus transform has a local continuity property (under some hypothesis) and  Proposition~\ref{prop_continuity_kernel} can in a certain way give us the local continuity of the frequency-focus transform. However none of the two results is strong enough to prove the existence of a Lipschitz constant that satisfies equation~\eqref{eq_lipchitz_constant}.

We plan to study this question in the future.
%In the linear case, both definition are equivalent, and even stronger they are equal which means that $C(F)=\rho(F)$ when $F$ is linear (the equality stand in $[0,+\infty]$ if $F$ is not necessarily continuous).
%
%It could be interesting to see if our two transforms have either their first or the second generalized norm finite. 
We also plan to investigate in which conditions there is a constant $c(M)$ such as
\begin{equation}
    \forall f_1,f_2\in L^2\ ,\quad  c(M)\|f_1-f_2\|_{L^2(\R)} \leqslant\|Mf_1 - Mf_2 \|_{L^2(\R^2)}\ .
\end{equation}
Such property would guarantee injectivity of the non-linear transform.

\smallskip
A main further goal will be to study the invertibility of such non-linear transforms. From our results, inverse transforms can be obtained if both the transform $M_f$ and the focus function $\sigma_f$ are known, but not in situations where only the transform $M_f$ is known. A first step would be to analyze in which conditions an approximate inverse can be obtained when an approximation of the focus function is available. This may open the door to iterative inversion methods.

\smallskip
Last but not least, we plan to head to concrete applications of this approach, in particular in the context of audio perception modelling. For that, we plan to investigate further focus functions that could be relevant in applications, starting from the toy models described in the appendix. Another key milestone will be the construction of a discrete theory, which will require appropriate sampling schemes (depending on the focus functions), probably in the spirit of~\cite{Balazs2011theory,Doerfler2015nonstationary}.
%Those results lead us to further purposes. The  first one would be to study the definitions of the focus function and to understand which properties of the signals we want to focus on. Another goal would be to estimate those focus function : for a given signal how to estimate with a certain precision $\eps$ the focus function associated to $f$. Having such estimate focus function we can study a linearized transform : instead of using the proper focus function of the signal, use the estimate focus function. Hence the obtained transform is linear and try to control the error with the main operator. Then consequence of the theorem \ref{theorem_norm_control_frequency_focus} and \ref{theo_bounding_time_marseillan_transform} is that for fixed focus functions ($\sigma_f^\nu$ and $\sigma_f^\tau$) the linearized transforms are injective which lead us to the idea of inversion of the linearized operator.